\documentclass[12pt]{article}
\usepackage[T1]{fontenc}
\usepackage[utf8]{inputenc}
\usepackage{lmodern}
\usepackage{textcomp}
\usepackage{enumitem}
\usepackage{geometry}
\geometry{margin=1in}
\begin{document}
\begin{center}
Answer Key - Biology - Graduate Level Assessment 21 \
\small (Answers are ordered exactly as in the question paper.)
\end{center}

\section*{Multiple Choice}
\begin{enumerate}[label=\arabic*.]
\item \textbf{Answer:} D
\\ \textit{Options:} A) O$_{2}$; B) CH 4 (Methane); C) CO 2; D) H$_{2}$O; E) N$_{2}$ (Nitrogen)
\\ \textit{Note:} The correct answer is H$_{2}$O because during photosynthesis, water molecules are split in a process called photolysis, releasing oxygen as a byproduct.
\item \textbf{Answer:} A
\\ \textit{Options:} A) The human genome has fully adapted to current environmental pressures; B) Modern medicine has eliminated the effects of genetic drift; C) The rate of genetic drift is directly proportional to technological advancements; D) Humans have reached a genetic equilibrium where drift is no longer possible; E) Humans have developed immunity to genetic drift
\\ \textit{Note:} The correct answer is based on the understanding that the human genome has undergone significant stabilization in response to environmental pressures over time, reducing the impact of random genetic drift compared to earlier periods when populations were more susceptible to such changes.
\item \textbf{Answer:} A
\\ \textit{Options:} A) Influenza, chickenpox, and measles; B) Multiple sclerosis, Parkinson's disease, and Alzheimer's disease; C) Arthritis, osteoporosis, and fibromyalgia; D) Hypertension, heart disease, and diabetes; E) Chronic obstructive pulmonary disease (COPD), cystic fibrosis, and pleural effusion
\\ \textit{Note:} Influenza, chickenpox, and measles are all viral infections that can significantly affect the respiratory tract, leading to various abnormalities such as inflammation, increased mucus production, and respiratory distress.
\item \textbf{Answer:} A
\\ \textit{Options:} A) the rate of mutations; B) the chance for variation in zygotes; C) the age of an organism; D) the potential for genetic disorders; E) the sequence of DNA bases
\\ \textit{Note:} Crossing-over during meiosis facilitates genetic recombination, which increases genetic diversity and allows scientists to assess mutation rates by analyzing the frequency and distribution of alleles in offspring.
\item \textbf{Answer:} A
\\ \textit{Options:} A) The parathyroid glands produce insulin to regulate blood sugar levels; B) The parathyroid glands regulate blood pressure; C) The parathyroid glands control breathing rate; D) The parathyroid glands secrete enzymes that break down proteins; E) The parathyroid glands control the balance of electrolytes in the body
\\ \textit{Note:} The parathyroid glands do not produce insulin; instead, they secrete parathyroid hormone (PTH), which regulates calcium levels in the blood, making the provided answer incorrect.
\end{enumerate}

\section*{True / False}
\begin{enumerate}[label=\arabic*.]
\setcounter{enumi}{5}
\item \textbf{Answer:} True
\\ \textit{Options:} A) True; B) False
\\ \textit{Note:} This testing method may provide a useful strategy for conducting HIV surveillance in possible co-infected TB patients at peripheral centres. Since there is no investment on infrastructure, it may be possible for paramedical health professionals to carry out the test, particularly in areas with low HIV endemicity.
\item \textbf{Answer:} True
\\ \textit{Options:} A) True; B) False
\\ \textit{Note:} Patients who received a metal-backed Onlay tibial component obtained better postoperative mechanical alignment compared to those who received all-polyethylene Inlay prostheses. The thicker overall construct of Onlay prostheses appears to be an important determinant of postoperative alignment. Considering their higher survivorship rates and improved postoperative mechanical alignment, Onlay prostheses should be the first option when performing medial UKR.
\item \textbf{Answer:} True
\\ \textit{Options:} A) True; B) False
\\ \textit{Note:} The findings showed that acupuncture of voice-related acupoints could bring about improvement in vocal function and healing of vocal fold lesions.
\item \textbf{Answer:} True
\\ \textit{Options:} A) True; B) False
\\ \textit{Note:} Our results indicate that elevated cTnI levels are associated with higher risk for inhospital mortality and complicated clinical course. Troponin I may play an important role for the risk assessment of patients with PE. The idea that an elevation in cTnI levels is a valuable parameter for the risk stratification of patients with PE needs to be examined in larger prospective studies.
\item \textbf{Answer:} True
\\ \textit{Options:} A) True; B) False
\\ \textit{Note:} These preliminary data support our hypothesis that recipient inflammation may affect RBC alloimmunization in humans; however, a more detailed understanding of the pathophysiologic association between inflammation and alloimmunization is required before definitive conclusions can be reached.
\end{enumerate}

\section*{Long Form Response}
\begin{enumerate}[label=\arabic*.]
\setcounter{enumi}{10}
\item \textbf{Answer:} Multiple alleles refer to the presence of three or more alternative forms of a gene that can occupy the same locus on a chromosome, influencing a single trait. This genetic phenomenon arises through mutations and evolutionary processes, leading to increased genetic diversity within a population. The existence of multiple alleles allows for a wider range of phenotypic expressions, as different combinations can result in varying traits. For example, in the ABO blood group system, three alleles ($I^A$, $I^B$, and i) determine blood type, showcasing how multiple alleles contribute to trait variation. This genetic complexity has significant implications for the adaptability and survival of populations, as it enhances their ability to respond to environmental changes.
\\ \textit{Note:} The answer correctly defines multiple alleles as alternative gene forms at a single locus, explains their origin through mutations and evolution, and illustrates their significance in increasing phenotypic diversity within populations, exemplified by the ABO blood group system.
\item \textbf{Answer:} Conjugation in paramecia is a process that typically facilitates genetic exchange and increases genetic diversity among populations. However, when paramecia descend from a single individual through repeated fission, all resulting individuals are genetically identical clones. This genetic uniformity precludes the necessity for conjugation, as there would be no genetic variation to exchange. Furthermore, the lack of genetic diversity could lead to reduced adaptability and resilience against environmental changes, emphasizing the importance of sexual reproduction in promoting genetic variability. Therefore, in a clonal population, conjugation would not occur due to the absence of genetic differences essential for this process.
\\ \textit{Note:} correct because it emphasizes that genetic identity among paramecia descended from a single individual eliminates the need for conjugation, as there is no genetic variation to exchange, which in turn limits genetic diversity and adaptability.
\end{enumerate}

\vfill
\noindent\textit{End of Answer Key}
\end{document}